% Generated draft: requires mathematical and source review.
\documentclass[11pt]{article}
\usepackage[margin=25mm]{geometry}
\usepackage{fontspec}
\setmainfont{Latin Modern Roman}
\usepackage{amsmath,amssymb,mathtools,unicode-math}
\setmathfont{Latin Modern Math}
\usepackage{enumitem,needspace,xurl,hyperref}
\setlength{\parindent}{0pt}
\setlength{\parskip}{5pt}
\setlength{\emergencystretch}{3em}
\newcommand{\N}{\mathbb{N}}
\newcommand{\Z}{\mathbb{Z}}
\newcommand{\Q}{\mathbb{Q}}
\newcommand{\R}{\mathbb{R}}
\newcommand{\C}{\mathbb{C}}
\newcommand{\F}{\mathbb{F}}
\newcommand{\A}{\mathbb{A}}
\newcommand{\PP}{\mathbb{P}}
\newcommand{\E}{\mathbb{E}}
\newcommand{\eps}{\varepsilon}
\newcommand{\dd}{\,\mathrm{d}}
\begin{document}
\section*{UnsolvedMath ID: 30003981}
\textbf{Source problem code:} catalogue entry
\subsection{Definitions and mathematical statement}
Let \(G\) be a group, let \(\mathbb F\) be a field, let \(V=\mathbb F^d\), and let \(\rho:G\to\operatorname{GL}(V)\) be a representation, where \(\operatorname{GL}(V)\) is the group of invertible \(\mathbb F\)-linear maps on \(V\). For \(z\in V\), define its orbit, interpreted literally as an unordered set, by
\[
\mathcal O_\rho(z)=\{\rho(g)z:g\in G\}.
\]
The supplied problem asks for a characterization of the linear subspaces \(S\subseteq V\) for which \(\mathcal O_\rho(z)\) uniquely determines almost every \(z\in S\), modulo the stated ``trivial scalar ambiguities.''

The source does not specify the field, the meaning of almost every, the allowed scalars, or the domain of competing vectors. To display these missing choices explicitly, suppose additionally that \(\mu_S\) is a specified measure on \(S\), that \(A\leq\mathbb F^\times\) is a specified group of allowed scalars, and that \(T\subseteq V\) is a specified competing domain containing \(S\). Write \(Az=\{az:a\in A\}\). A measure-based interpretation of the requested recovery property is that there exists a \(\mu_S\)-null measurable set \(N\subseteq S\) such that
\[
\forall z\in S\setminus N\quad\forall w\in T,\qquad
\mathcal O_\rho(w)=\mathcal O_\rho(z)
\ \Longrightarrow\ w\in Az.
\]
Under these additional choices, the task would be to characterize exactly the subspaces satisfying this property. This is a conditional formulation, not a resolution of the source's missing specifications. In particular, the source does not choose between \(T=S\) and \(T=V\), does not identify \(A\), and does not state whether the data are invariant under all allowed scalars. It also does not determine whether ``almost every'' means measure-theoretic almost everywhere or an algebraic notion of genericity.

There is a further unresolved issue concerning the observation itself: a full unordered orbit, a family indexed by \(G\), and phase-retrieval measurements are different data. An unordered orbit carries ambiguities arising from the group action that need not be scalar. The source does not explain how those ambiguities are to be treated. Consequently, the intended precise recovery problem requires clarification.
\subsection{Short English statement}
For a given group representation, characterize the subspaces on which almost every vector can be recovered from its orbit up to specified scalar ambiguities. The supplied record leaves the field, genericity convention, competing vectors, allowed scalars, and exact observed data unresolved, so the displayed quantified statement is conditional on those choices.
\subsection{Sources}
\begin{itemize}
\item Supplied local UnsolvedMath registry record, permanent ID \(30003981\), titled ``Generic Recovery from Group Orbits on Subspaces,'' in the Oberwolfach Reports Open Problems collection; exact source problem code \texttt{catalogue entry}. No individual proposer is named. The record gives proposed year \(2018\).
\item Supplied source citation: \emph{Mini-Workshop: Algebraic, Geometric, and Combinatorial Methods in Frame Theory (2019)}, Oberwolfach Reports, pp.~2769--2803, DOI \texttt{10.4171/owr/2018/46}. \url{https://doi.org/10.4171/owr/2018/46}. The supplied background describes these reports as documenting mathematical workshops and their open-problem discussions. The citation's \(2019\) designation and the record's \(2018\) year are retained without reconciliation.
\item Supplied literature source labeled ``Algebraic phase retrieval'': \url{https://www.research.ed.ac.uk/files/58341345/1312.0158.pd.pdf}.
\item Supplied literature source labeled ``Phase retrieval in dimension four'': \url{https://arxiv.org/abs/2607.27719}.
\item Literature assessment dated exactly \texttt{2026-08-21}: the supplied curation labels the problem open and states that the real case is settled, complex generic sufficiency is known, and counterexamples to the \(4d-4\) lower bound together with unresolved dimensions leave the general minimum open. These are supplied dated curation claims, not independently verified conclusions. Their connection to the stated orbit-recovery problem is unspecified. The supplied background itself requests independent verification of time-sensitive status and cited claims.
\item The registry labels the statement \texttt{corrected\_verified} and says that notation, typography, and sentence structure were repaired while preserving mathematical content. This is supplied editorial metadata, not independent mathematical verification of this draft or its interpretation.
\end{itemize}
\end{document}
